\begin{textsample}{POS Dim 2 – sc – Score 19.00 – sc/sc\_em\_1.txt}  \label{ex:f2_pos_020}
APRESENTAÇÃO > “ Eu quero desaprender para aprender de novo. > Raspar as tintas com que me pintaram. > Desencaixotar emoções, recuperar sentidos ”. > ( Rubem Alves ) Partindo dos dizeres de Rubem Alves, apresentamos este \textbf{Currículo} Base do Ensino Médio do Território Catarinense, indicando que ele exige o repensar dos modos e das práticas do fazer escolar, numa constante dinâmica de desaprender para aprender.
Este documento, que agora se apresenta, é uma potencial ferramenta norteadora do trabalho pedagógico, no qual os múltiplos saberes se relacionam a partir da organização e da teorização do conhecimento em áreas, com vistas a colaborar com a formação do estudante em sua totalidade e de forma conectada com sua realidade.
Espera -se que esta proposta de organização curricular ressoe em um ensino significativo e atrativo, permitindo aos estudantes serem protagonistas de suas histórias no planejar seus percursos escolares e suas vidas com maior consistência, desenvolvendo -se integralmente.
A construção deste \textbf{currículo}, alicerçado pela Base Nacional Comum Curricular ( BNCC ), foi pensada no sentido de garantir a mobilização necessária para a participação democrática.
Nessa perspectiva, os trabalhos ocorreram \textbf{alinhados} com a União Nacional dos Conselhos \textbf{Municipais} de Educação ( Uncme/SC ), a União Nacional dos \textbf{Dirigentes} \textbf{Municipais} de Educação ( Undime/SC ), a Federação Catarinense de Municípios ( Fecam/SC ), o Conselho Estadual de Educação ( CEE/SC ) e a Secretaria de Estado da Educação ( SED/SC ).
O resultado deste processo se deu a partir de um intenso trabalho coletivo, que contou com a participação de mais de 300 \textbf{profissionais} da Rede Estadual de Ensino, atuantes nas diversas esferas da educação.
Esta construção democrática reflete a busca contínua por um sólido percurso de construção da Proposta Curricular do Estado, a qual imprime, como neste documento, a multiplicidade das vozes de \textbf{profissionais} da educação.
Vale registrar que este documento - \textbf{Currículo} Base do Ensino Médio do Território Catarinense - complementa o \textbf{Currículo} Base da Educação Infantil e do Ensino Fundamental do Território Catarinense.
Logo, deve -se compreendê -lo como um continuum processo de discussões, a partir do documento já \textbf{homologado}, uma vez que este se constitui como complementaridade daquele.
Sendo assim, a educação catarinense passa a contar, a partir de agora, com um documento completo que norteará a educação básica para todo o sistema de ensino, que servirá como um material de apoio para os educadores de todas as etapas do ensino básico e suas modalidades, passando a trabalhar na perspectiva das Áreas do Conhecimento, conectando teoria e prática.
Para contextualizar este processo, registramos que a elaboração deste \textbf{Currículo} Base do Ensino Médio foi iniciada em abril de 2019, através do Programa de Apoio à Implementação da Base Nacional Comum Curricular ( ProBNCC ).
Envolveu -se neste trabalho uma equipe de vinte e cinco \textbf{profissionais} da Rede Estadual de Ensino de Santa Catarina, apoiada por uma consultora geral de \textbf{currículo}.
Esta equipe realizou o estudo dos principais documentos de referência - a Proposta Curricular de SC ; o \textbf{Currículo} Base da Educação Infantil e do Ensino Fundamental do Território Catarinense ; a Base Nacional Comum Curricular e as \textbf{Diretrizes} Curriculares Nacionais para o Ensino Médio ( 2018 ) -, possibilitando, assim, a construção da primeira versão do documento, denominado"Marco Zero".
Essa proposição inicial foi submetida a consulta pública, que recebeu 2.120 contribuições válidas nas diferentes áreas do conhecimento.
Após este movimento inicial, foram integrados ao processo, por meio de um edital de seleção que dispunha vagas para professores, gestores e \textbf{profissionais} das Coordenadorias Regionais de Educação de todas as mesorregiões do estado, um total de 254 \textbf{profissionais} da educação efetivos da Rede Estadual de Ensino, os quais foram convidados a atuar como professores elaboradores/colaboradores/formadores, dando continuidade à escrita do documento.
Foi programado para participarem, de forma presencial, de um ciclo de três seminários, quando iriam contribuir para o aprofundamento do documento.
Porém, estes seminários foram reprogramados por conta da pandemia da Covid-19, elaborando -se um novo cronograma para ocorrer de forma remota ( online ).
Assim, a partir do mês de junho de 2020, os trabalhos foram retomados.
Além disso, mobilizaram -se uma equipe de \textbf{profissionais} das diferentes \textbf{diretorias} da Secretaria de Estado da Educação e uma equipe de consultores de cada uma das áreas de conhecimento para a estruturação e concretização do \textbf{currículo}. \textbf{Alinhados} a esse processo, 363 educadores das 120 escolas-piloto do Novo Ensino Médio, no âmbito do ProNem ( Programa de Apoio a Núcleos Emergentes ), apoiados pela equipe \textbf{técnica} da SED, por especialistas externos, com base nas experiências vivenciadas no processo de implementação do Novo Ensino Médio e nos parâmetros da BNCC, elaboraram o portfólio dos Componentes Curriculares \textbf{Eletivos} que integram este documento.
Uma vez que o trabalho estava sendo realizado de forma remota, muitos encontros online foram estabelecidos, partindo de um cronograma geral e também por iniciativa dos grupos de trabalho, que, de forma uníssona, tinham como objetivo elaborar um novo \textbf{currículo}, capaz de produzir novos sentidos e significados aos estudantes, bem como o de favorecer o protagonismo juvenil na realização das escolhas dos \textbf{itinerários} \textbf{formativos} em \textbf{conformidade} com seus projetos de vida.
Nessa perspectiva, acredita -se que este documento sirva para romper com anos de desinteresse de grande parte dos estudantes, fator que contribui significativamente para os altos índices de abandono escolar e a baixa proficiência, que \textbf{demarcam}, historicamente, a realidade do Ensino Médio no Brasil.
Espera -se, também, que esta proposta contribua para fortalecer o sentimento de pertencimento de estudantes e professores para com a escola, fator determinante para uma educação de \textbf{qualidade} e com melhores resultados.
Por fim, conclui -se que este documento, ao mesmo tempo que representa o nosso ponto de chegada, representa também o grande ponto de partida, pois irá referenciar o trabalho educacional por longos anos, tornando -se, assim, um documento vivo, que produza sentidos e significados para os estudantes do Ensino Médio de Santa Catarina.
Assim, convidam -se todos os educadores a seguir juntos no desafio de implementar este \textbf{currículo} na prática, nas escolas, que é onde a educação acontece, na perspectiva de que se concretize, neste estado, uma educação pública de \textbf{qualidade}!
Finalmente, não poderíamos deixar de mencionar a atipicidade do momento de construção deste importante documento, tecido ao longo de um ano marcado por uma pandemia mundial.
Contudo, mesmo diante de todos os desafios e \textbf{adversidades} enfrentadas, saímos vitoriosos : entregamos um \textbf{currículo} que reflete a \textbf{excelência} dos \textbf{profissionais} da Rede Estadual de Ensino.
Luiz Fernando Cardoso \textbf{Secretário} de Estado da Educação

% matched lemmas: adversidade, alinhar, conformidade, currículo, demarcar, diretoria, diretriz, dirigente, eletivo, excelência, formativo, homologar, itinerário, municipal, profissional, qualidade, secretário, técnico
\end{textsample}
